\chapter{Introducere in Python}
\label{chapter:python-intro}

O sa presupun ca stii deja cum se parcurg vectori si matrici, ai facut
asta de multe ori la C++. Fara presupunerea asta ne-am duce prea mult in
spate cu recapitularea, si nu la asta ne intereseaza timpul. Ce facem aici
e sa traducem ce stii deja in sintaxa lui Python, cat mai repede, ca sa
ajungem la partea interesanta.

\section{De ce Python, nu C++}

Intrebarea e corecta. Ai invatat C++ pentru olimpiada pentru ca e rapid,
apropiat de masina, te lasa sa controlezi fiecare octet. De ce am schimba
asta acum?

Raspunsul scurt: nu il schimbam, il ascundem. Aproape tot ce faci in
NumPy sau PyTorch ruleaza, pe dedesubt, in cod C sau C++ compilat, ai
vazut deja asta in \hyperref[chapter:numpy-pandas]{capitolul urmator}
cand comparam \pycode{v.mean()} cu bucla scrisa de mana. Python nu
inlocuieste C++, il apeleaza.

Ce castigi renuntand la C++ ca limbaj de zi cu zi:

\begin{itemize}
  \item scrii mai putin cod pentru aceeasi idee, deci iterezi mai repede;
  \item ai acces la un ecosistem urias de librarii facute exact pentru
    ML (NumPy, Pandas, scikit-learn, PyTorch);
  \item gresesti mai rar la lucruri marunte (memorie, tipuri), pentru ca
    Python le gestioneaza singur.
\end{itemize}

Ce pierzi: viteza bruta pe cod scris de mana. Dar in ML aproape niciodata
nu scrii tu bucla critica, o scrie altcineva o data, in C, si tu o apelezi
de mii de ori. De-aia insist sa intelegi ce e in spate, nu ca sa scrii tu
C prin Python, ci ca sa stii cand o functie e doar sirop peste ceva
simplu si cand chiar face treaba grea pe care n-ai vrea sa o repeti
singur.

\section{Cum ar trebui sa scriem Python}

Cateva reguli care difera fata de C++, si pe care merita sa le ai in cap
de la inceput:

\begin{itemize}
  \item Blocurile de cod sunt definite prin indentare, nu prin acolade.
    Daca indentarea e gresita, codul e gresit.
  \item Nu declari tipul unei variabile. \pycode{x = 5} e de ajuns,
    Python afla singur ca \pycode{x} e intreg.
  \item Nu pui punct si virgula la finalul liniei.
  \item Numele de variabile si functii se scriu \pycode{snake_case}, nu
    \pycode{camelCase}.
  \item Daca te trezesti scriind o bucla explicita peste o lista mare de
    date, intreaba-te intai daca nu exista deja o functie care face
    asta, o sa revenim serios la ideea asta in capitolul urmator.
\end{itemize}

\section{Sintaxa pe repede inainte}

Trecem prin echivalentele directe ale lucrurilor pe care le stii deja din
C++. Nu explic ce e o bucla sau ce e o functie, doar cum arata in Python.

\begin{minted}{python}
# variabile, fara tip declarat
n = 5
nume = "Alex"
gasit = True

# if / else
if n > 0:
    print("pozitiv")
elif n == 0:
    print("zero")
else:
    print("negativ")

# bucla, echivalentul lui for (int i = 0; i < n; i++)
for i in range(n):
    print(i)

# bucla, echivalentul lui for (int x : v)
v = [1, 2, 3]
for x in v:
    print(x)

# functie
def patrat(x):
    return x * x
\end{minted}

\pycode{list}-ul e echivalentul lui \pycode{vector} din C++, doar ca
poate tine orice tip de date amestecat si isi schimba singur dimensiunea:

\begin{minted}{python}
v = []
v.append(10)
v.append(20)
print(len(v))   # 2
print(v[0])     # 10
\end{minted}

Operatorul \texttt{\%} (rest) si operatorul \pycode{//} (cat la
impartirea intreaga) fac exact ce fac \texttt{\%} si \texttt{/} intre
intregi in C++, doar ca in Python \pycode{/} intre intregi da mereu
rezultat zecimal, de-aia exista \pycode{//} separat pentru impartirea
intreaga.

O sa vezi in capitolul urmator ca, pentru date numerice, NumPy inlocuieste
\pycode{list}-ul cu ceva mai apropiat, ca structura de memorie, de
\pycode{vector}-ul din C++. Dar pana acolo, \pycode{list}-ul e suficient.
